\section*{Glosario y abreviaturas}
\addcontentsline{toc}{section}{Anexo A: Glosario y abreviaturas}

\subsection*{Glosario de términos}
\begin{itemize}
    \item \textbf{Módulo}: Componente de software que añade funcionalidades específicas a un sistema mayor. En este proyecto, se refiere a los módulos de Gestión de Tiempos de Jornada y Gestión de Evidencias SINAES.
    \item \textbf{Tiempo de Jornada}: Unidad de medida interna de la Universidad Nacional que representa un valor similar a créditos o recursos en horas, utilizada para asignaciones docentes, proyectos institucionales y presupuesto de personal académico.
    \item \textbf{Horas Contacto}: Horas de clase presenciales o virtuales asociadas a un curso, fundamentales para el cálculo del tiempo de jornada según la normativa universitaria.
    \item \textbf{Malla Curricular}: Estructura oficial del plan de estudios de una carrera, que detalla la secuencia de cursos y sus respectivas horas contacto y créditos.
    \item \textbf{Curso de Repitencia}: Curso que un estudiante debe cursar nuevamente y que puede implicar una asignación adicional de tiempo de jornada para el profesor.
    \item \textbf{Proveedor Externo (Tiempos)}: Entidad o convenio externo a la UNA que otorga tiempos de jornada o recursos equivalentes para actividades específicas.
    \item \textbf{Evidencia}: Documento o registro formal que sirve como prueba del cumplimiento de un criterio o estándar de calidad, especialmente en acreditación.
    \item \textbf{Dimensión SINAES}: Nivel superior de la estructura de acreditación del SINAES, que agrupa varios componentes relacionados.
    \item \textbf{Componente SINAES}: Nivel intermedio de la estructura SINAES, que desglosa una dimensión en áreas más específicas.
    \item \textbf{Criterio SINAES}: Nivel más granular de la estructura SINAES, define un estándar específico que se verifica con evidencias.
    \item \textbf{Autenticación}: Proceso de verificar la identidad de un usuario que intenta acceder al sistema.
    \item \textbf{Autorización}: Proceso de determinar qué acciones o recursos puede acceder un usuario una vez autenticado, según sus roles.
    \item \textbf{Cifrado AES-256}: Estándar de cifrado con clave de 256 bits para proteger datos sensibles.
    \item \textbf{Metodología Ágil}: Enfoque de gestión basado en entrega iterativa e incremental, colaboración y adaptación al cambio.
    \item \textbf{Scrum}: Marco ágil con sprints, dailies, reviews y retrospectivas para gestionar proyectos.
    \item \textbf{API}: Interfaz de programación que permite la comunicación entre componentes o servicios.
    \item \textbf{ORM}: Mapeo Objeto–Relacional para interactuar con la base de datos mediante objetos.
    \item \textbf{Desarrollo iterativo}: Ciclos repetitivos de planificar–desarrollar–probar con retroalimentación continua.
\end{itemize}

\subsection*{Abreviaturas}
\begin{itemize}
    \item \textbf{UNA}: Universidad Nacional de Costa Rica
    \item \textbf{SINAES}: Sistema Nacional de Acreditación de la Educación Superior
    \item \textbf{SRS}: Software Requirements Specification (Especificación de Requerimientos de Software)
    \item \textbf{EDT / WBS}: Estructura de Desglose del Trabajo (Work Breakdown Structure)
    \item \textbf{API}: Application Programming Interface
    \item \textbf{ORM}: Object–Relational Mapping
    \item \textbf{JWT}: JSON Web Token
    \item \textbf{RBAC}: Role–Based Access Control
    \item \textbf{WCAG}: Web Content Accessibility Guidelines
    \item \textbf{UI}: User Interface
    \item \textbf{UX}: User Experience
    \item \textbf{PDF}: Portable Document Format
    \item \textbf{XLSX}: Formato de hoja de cálculo de Microsoft Excel
    \item \textbf{PNG}: Portable Network Graphics
    \item \textbf{DOCX}: Formato de documentos de Microsoft Word
    \item \textbf{CI/CD}: Continuous Integration / Continuous Delivery
    \item \textbf{DDD}: Domain–Driven Design
    \item \textbf{SME}: Subject Matter Expert
    \item \textbf{PMO}: Project Management Office
\end{itemize}
